% =============================================================================
%  Section 2b — Threat Model
% =============================================================================
\section{Threat Model}
\label{sec:threat}

This section makes explicit the adversary's capabilities and the boundary
between what Neuron's coupled governance guarantees and what lies outside
its scope. The boundary is restated in Section~\ref{sec:eval-setup} as
\emph{methodological boundaries} of the evaluation; here we establish the
structural assumptions on which the security claim rests. We follow the
convention of separating \emph{attacker profile} (who the adversary is and
what they control), \emph{attacker capabilities} (what they can do at
runtime), and \emph{out-of-scope} (what the runtime does not defend
against), as is standard in OS-security papers
~\cite{klein2009sel4,mercer2012capsicum}.

% -----------------------------------------------------------------------------
\subsection{Attacker Profile}
\label{sec:threat-profile}

We assume a \emph{rational, possibly colluding} adversary who controls at
least one tenant in the multi-tenant Neuron runtime. Concretely, the
adversary is the operator of one or more \emph{malicious tenant accounts}
that share the same runtime instance as benign tenants; the agents
instantiated under these accounts execute arbitrary, attacker-supplied
prompts and tool calls. We do not assume the adversary has kernel
privileges, root inside a sandbox, or any out-of-band control over the
host OS: the attack surface is confined to the abstractions the runtime
exposes to tenant agents (the VFS, the EventBus, the syscall interface,
the spawn path, and the approval channel).

We model three concrete attacker profiles, in increasing capability:

\begin{description}
    \item[P1. Noisy co-tenant.] The adversary runs one or more agents
    under a tenant account and floods the shared substrates with
    high-frequency \IO{} and broadcasts. The goal is to starve the
    governance layer and degrade co-resident tenants. The adversary
cannot bypass the runtime API; every operation traverses
\texttt{VfsManager} and \texttt{EventBus}.

    \item[P2. Privilege-seeking sub-agent.] The adversary's agent has
    been deliberately downgraded (for example, a reviewer restricted to
    a \texttt{*:deny} profile) and attempts to recover privileges
    \emph{through the spawn path}---by spawning a sub-agent that
    inherits a fresh default profile---or by phrasing a destructive-tool
    request in natural language to slip past depth-oblivious approval.

    \item[P3. Cross-tenant infiltrator.] The adversary combines P1 and
    P2: under cover of resource contention, the spawned sub-agent
    attempts a cross-tenant VFS access while the permission layer is
impaired. This is the strongest profile the runtime must defend
against and is the configuration exercised by Scenario~6
    (Section~\ref{sec:eval-cross}).
\end{description}

% -----------------------------------------------------------------------------
\subsection{Attacker Capabilities}
\label{sec:threat-cap}

Within the boundary above, we grant the adversary the following
runtime capabilities, each of which corresponds to an attack surface
characterized in Section~\ref{sec:bg-attack} and closed by a mechanism
of Section~\ref{sec:arch}:

\begin{itemize}
    \item \textbf{Issue arbitrary prompts and tool calls.} The
    adversary's agent may submit any prompt, request any tool exposed
    by the runtime, and phrase a privilege-escalation request in any
    natural language. The runtime's defense is structural
    (\texttt{EscalationPolicy} adjudicates depth and tool class, not
    language) rather than prompt-level.

    \item \textbf{Spawn sub-agents at runtime.} The adversary's agent
    may invoke \texttt{AgentTool} to create sub-agents, producing a
    spawn tree of unbounded breadth up to the configured depth cap.
    Sub-agents inherit the attacker's tenant identity.

    \item \textbf{Consume the shared substrates.} The adversary's agents
    may issue arbitrarily many VFS reads/writes and EventBus
    broadcasts, subject only to the runtime's rate limiters. The
    attacker is assumed to know the per-tenant limits (these are
    documented) and may attempt to operate just below them
    (\emph{low-and-slow} attacks; see Section~\ref{sec:eval-adaptive}).

    \item \textbf{Observe their own audit trail.} The adversary can
    read the \texttt{UnifiedAuditLog} entries their own actions
    produce; they cannot read, tamper with, or selectively delete
    other tenants' entries. The audit chain is append-only at the file
    system level and is not exposed for in-place mutation through any
    runtime API.

    \item \textbf{Time their operations.} The adversary can observe
    wall-clock latency of their own calls and use this as a side
    channel to infer the state of the shared lock or the rate limiter.
    This is the standard timing side channel of any shared-resource
    system; Section~\ref{sec:lim-failure} discusses its limits.
\end{itemize}

% -----------------------------------------------------------------------------
\subsection{Out-of-Scope}
\label{sec:threat-out}

The following are explicitly \emph{not} defended against, and the
evaluation does not claim otherwise:

\begin{itemize}
    \item \textbf{Prompt injection at the model layer.} The runtime
    does not inspect prompt content for jailbreaks or indirect prompt
    injection; that is the responsibility of upstream guardrails
    (Section~\ref{sec:rw-sec}). What the runtime guarantees is that
    \emph{once a malicious escalation request reaches the policy gate},
    it is denied regardless of phrasing---a prevention/enforcement
    separation restated in Section~\ref{sec:eval-setup}.

    \item \textbf{Kernel or host compromise.} If the adversary has
    kernel privileges, host root, or the ability to mutate the JVM or
    the on-disk audit log, the security claim is void. This is the
    standard trust-the-TCB assumption; seL4~\cite{klein2009sel4} and
    Capsicum~\cite{mercer2012capsicum} make the same boundary
    explicit.

    \item \textbf{Tool-execution-to-JVM reverse attack.} A natural
    concern is that an agent that legitimately obtains
    \texttt{bash}---which Section~\ref{sec:bg-attack} acknowledges is
    nearly a defining capability of an LLM agent---could use that shell
    to tamper with the JVM process hosting the runtime (via
    \texttt{/proc}, \texttt{ptrace}, signals, or filesystem writes to
    the JVM's own state). We close this boundary \emph{structurally
    rather than by policy}: every tool invocation in Neuron traverses
    the \texttt{SandboxProvider} abstraction
    (Section~\ref{sec:arch-sandbox}), which offers two isolation
    regimes. \emph{In-process WASM isolation} executes compiled
    WebAssembly within the GraalVM polyglot engine with a hard memory
    ceiling, and sandboxed code cannot reach host resources
    directly---it issues semantic syscalls bound to runtime-mediated
    host functions, so there is no \texttt{bash} process at all and no
    path to JVM internals. \emph{Out-of-process container isolation}
    executes untrusted code in a disposable Docker container with
    \texttt{--cap-drop ALL} (no Linux capabilities),
    \texttt{--security-opt no-new-privileges} (no \texttt{setuid}
    escalation), \texttt{--security-opt seccomp=<profile>} that
    explicitly blocks container-escape syscalls (\texttt{unshare},
    \texttt{pivot\_root}, \texttt{mount}, \texttt{ptrace}),
    \texttt{--read-only} root filesystem, \texttt{--network none}
    (air-gapped), and \texttt{--pids-limit 64}. The JVM process runs
    on the host, outside the container's mount namespace and PID
    namespace; the agent's \texttt{bash} process cannot see the JVM
    process, cannot \texttt{ptrace} it (seccomp blocks
    \texttt{ptrace} and the container lacks \texttt{CAP\_SYS\_PTRACE}),
    cannot write to the JVM's filesystem (read-only rootfs, no shared
    mounts), and cannot signal it (different PID namespace). The
    agent's bash is therefore confined to the container boundary, and
    the JVM is structurally outside its reach. We note that the
    in-process WASM path is the default for whitelisted tools, while
    the container path is mandatory for any tool not on the whitelist
    (including any newly registered shell-like tool); this makes
    ``trust boundary = container boundary'' the default for
    non-whitelisted execution, rather than a user opt-in.

    \item \textbf{Side channels outside the runtime's view.} Cache-line
    timing, power analysis, and similar physical side channels are out
    of scope. The timing side channel on the shared VFS lock is in
    scope but is bounded by the rate limiter
    (Section~\ref{sec:eval-adaptive}).

    \item \textbf{Adversarial model weights or training-time
    poisoning.} We assume the underlying LLM weights are not
    adversary-controlled. A model whose weights have been poisoned to
    emit specific escalation requests is upstream of the runtime's
    enforcement boundary; the runtime still blocks the resulting
    request, but the prevention layer must be addressed elsewhere.

    \item \textbf{Inter-tenant collusion via out-of-band channels.}
    Two adversaries on different tenants may coordinate out-of-band
    (e.g.\ via a shared clock), but each operates under the
    per-tenant limits; the runtime's per-tenant accounting is
    unaffected by whether the tenants are coordinating.
\end{itemize}

% -----------------------------------------------------------------------------
\subsection{Boundary against the Evaluation}
\label{sec:threat-eval}

The boundary above has two direct consequences for the evaluation
that warrant explicit statement, because they bound the claims a
reviewer should accept.

\paragraph{Prevention vs.\ enforcement.}
Scenario~6 (Section~\ref{sec:eval-cross}) measures
\emph{runtime-level enforcement}: given that a malicious escalation
request has reached the policy gate, does the runtime deny it? It does
not measure the upstream probability that a compromised agent would
emit such a request; that probability is the domain of prompt-injection
research. The supplementary real-LLM experiment
(Section~\ref{sec:eval-supplement}) closes part of this gap by showing
that real LLMs do spontaneously emit such requests under adversarial
goals, but it is not a prevention-layer evaluation.

\paragraph{Simulated vs.\ real resource rates.}
The resource-layer benchmark (Scenario~4,
Section~\ref{sec:eval-resource}) drives the cgroup at simulated
generation rates ($U[15,55)$\,ms inter-increment jitter) rather than
by real LLM inference. This is because the measurement target is the
\texttt{preCheckAndReserve} gateway and the OOM-trigger latency, both
of which are functions of the cgroup accounting path and are
independent of whether the tokens are produced by a real model or by
a deterministic emitter. The supplementary end-to-end experiment
(Section~\ref{sec:eval-e2e}) closes the realism gap by replaying the
same attack with a real LLM in the loop, quantifying the absolute
latency cost in the regime where LLM inference dominates.

\paragraph{Marginal security value when upstream guardrails fail.}
A reviewer may ask: if the dominant real-world attack vector is prompt
injection, and Neuron explicitly does not defend against it, what is the
marginal value of a runtime-level enforcement layer? We argue the value
is structural and is not refuted by the prevalence of prompt injection.
\emph{First}, prompt-injection defenses are best-effort and known to be
defeated by paraphrase, encoding, and indirect injection through
retrieved content~\cite{debenedetti2024agentdojo}; treating them as a
single layer of defense places the entire security posture on a
probabilistic mechanism. A runtime enforcement layer that is
\emph{deterministic} (the Scenario~6 policy is a pure function of depth
and tool class) provides defense-in-depth: the worst case for upstream
guardrails becomes ``the malicious request reaches the policy gate,''
not ``the malicious request executes.'' \emph{Second}, the coupled
governance property is independent of how the request arrived: the
starvation and joint-decision experiments
(Sections~\ref{sec:eval-starvation}, \ref{sec:eval-joint}) demonstrate
that even under active resource contention---a regime in which upstream
guardrails themselves degrade---the runtime still denies the request.
\emph{Third}, the spawn-path closure (privilege non-increase) catches
the case where a prompt-injected parent agent attempts to spawn a
sub-agent to recover privileges it has already lost; no upstream
guardrail that operates only on prompt content can detect this, because
the spawn operation is structurally invisible at the prompt layer. The
supplementary real-LLM experiment (Section~\ref{sec:eval-supplement})
provides direct evidence: real commercial LLMs, given adversarial goals
but \emph{no} specific instruction to attack the spawn path,
spontaneously emit escalation requests that the runtime intercepts. We
therefore frame the contribution as a \emph{last-line enforcement
guarantee}, not a replacement for upstream prevention; the two layers
are complementary, and the absence of prevention does not void the
enforcement guarantee.

These two boundaries are the structural expression of the
prevention/enforcement and simulation/realism distinctions that
Section~\ref{sec:eval-setup} restates as methodological boundaries of
the experimental design.
